\subsection{A geometric vanishing theorem and necessary fibre defects}

\begin{statement}{C4a}{Proved}{A geometric replacement on an explicit fibre class}
Let $f:Y\to\PP^1$ be a connected-fibre holomorphic map from a compact connected complex threefold with $b_1(Y)=b_2(Y)=0$, and smooth general fibre a complex two-torus. Suppose every full fibre has one of the following forms: a reduced irreducible normal-crossing surface; or a smooth torus family modulo a finite free cyclic action over a cyclic base cover, so that its reduction is smooth and its multiplicity is the covering order. Let
\[
 0\ne\gamma\in H^0(\PP^1,R^1f_*\Z)
\]
be a generator, and suppose that, for every singular reduced full fibre $S$, its pullback to $H^1(\widetilde S,\Q)$ under normalization is nonzero. Then a dense open subset of every component of $\Pic(Y)$ satisfies
\[
 H^q(Y,T_Y\otimes L)=0\quad(0\le q\le3),\qquad
 \chi(Y,T_Y\otimes L)=0,\qquad c_3(Y)=0.                \tag{C.15}
\]
The generator exists under the stated rational Betti-number assumptions. Its normalization condition is a condition on an explicit integral cohomology class, with no tangent-cohomology vanishing assumed.
\end{statement}
\begin{proof}
The integral Leray sequence gives an injection
\[
H^0(\PP^1,R^1f_*\Z)\longrightarrow H^2(\PP^1,\Z)=\Z
\]
whose cokernel injects into the finite group $H^2(Y,\Z)$. Thus its source is infinite cyclic. The degree-one specialization argument of Part D identifies its stalk at every fibre with a nonzero integral class: a class in the preimage of a small disc that vanishes on a nearby fibre comes from the meridian of the punctured disc; a normal disc to any irreducible component makes a nonzero multiple of that meridian nullhomotopic, forcing the class to vanish. The real-to-holomorphic cohomology map on a compact connected complex space is injective on real classes, as proved there by the pluriharmonic maximum principle. Coherent Leray therefore supplies $\beta\in H^1(Y,\OO_Y)$ whose restriction to every full fibre is the image of $\gamma$.

Fix a real irrational number $r$ and put $M_r=\exp(r\beta)$, with the exponential normalized by $2\pi i$. On each fibre this line is represented by the unitary character $\exp(2\pi ir\gamma)$. It is nontrivial, and on a singular fibre its normalization pullback is nontrivial by the hypothesis on $\gamma$. To see this directly, triviality of an $n$th tensor power would place $nr\gamma$ in the integral lattice after applying $H^1(\R)\to H^1(\OO)$; real injectivity makes this an equality in $H^1(\R)$, which is impossible for an irrational multiple of a nonzero integral vector.

The same construction works in every Picard component. Each component contains a finite-order line bundle $P$: for a representative $P_0$, a positive power $P_0^m$ lies in the complex vector group $\Pic^0(Y)$, where division by $m$ is possible. Dividing $P_0$ by the resulting line makes its $m$th power trivial. Replacing $M_r$ by $M_r\otimes P$ retains nontrivial unitary restriction on each fibre and on each specified normalization. Indeed an $m$th power eliminates $P$ and reduces a purported trivialization to the preceding irrationality contradiction.

Set $L=K_Y\otimes(M_r\otimes P)^{-1}$. On every full fibre $F$ the Serre-dual bundle for degree two is
\[
 (T_Y\otimes L)|_F^*\otimes\omega_F
       =\Omega_Y^1|_F\otimes(M_r\otimes P)|_F,
\]
because $F$ is a Cartier fibre and $\omega_F=K_Y|_F$. At a smooth torus fibre both the conormal line and the intrinsic cotangent bundle are trivial; nontrivial unitary degree-zero lines on the torus have zero cohomology. The ambient space of sections is therefore zero. At a reduced irreducible normal-crossing fibre, Statement A.4.5 proves the torsion-free cotangent vanishing for every nontrivial unitary line. The normal-crossing kernel in A.1.2 is the normalization pushforward, and its sections vanish because the pulled-back unitary line is nontrivial. A holomorphic section of a unitary line on a compact connected complex manifold has constant norm by the maximum principle and is either zero or a trivialization. Thus the ambient group vanishes there as well.

On a cyclic-cover neighbourhood of a smooth multiple fibre, the covering map to $Y$ is finite etale. The pulled-back line restricts nontrivially to the covering central torus: finite etale pullback is injective on rational degree-one cohomology by its trace, and $r$ is irrational. The scheme-theoretic inverse image of the full fibre is the corresponding thickening of this torus. Each layer in its ideal filtration is the same ambient cotangent restriction tensored with this nontrivial line, since the covering base coordinate trivializes the conormal. Every layer has zero cohomology. The exact filtration sequences therefore give zero sections on the full thickening. Taking finite-group invariants, an exact functor, proves the required ambient vanishing on the original full fibre.

Top-degree coherent base change now gives $R^2f_*(T_Y\otimes L)=0$. For sufficiently large integer $n$, Serre vanishing on the base gives
\[
H^1\bigl(\PP^1,R^1f_*(T_Y\otimes L)\otimes\OO(n)\bigr)=0.
\]
Projection formula and the two-column Leray sequence imply
$H^2(Y,T_Y\otimes L\otimes f^*\OO(n))=0$.
This is an actual twist in the same Picard component up to the fixed finite-order class of $f^*\OO(n)$; choosing $n$ divisible by the finite order of its Chern class keeps the component unchanged. Semicontinuity in the explicit Picard family of C.1 makes the degree-two zero locus dense open in that component. Intersect it with the dense open endpoint-vanishing loci from C.1. On this intersection $\chi(T_Y\otimes L)=-h^1\le0$.

Statement A.4.6 gives Euler number zero or two for each reduced irreducible normal-crossing fibre; the positive case has rational normalization. Our normalization hypothesis excludes that case, since rational smooth surfaces have $H^1(-,\Q)=0$. Smooth torus fibres and the smooth finite etale quotients in the second fibre class have Euler number zero. Therefore C2a gives $c_3(Y)=0$. The index formula in C.1 gives $\chi=0$, forcing $h^1=0$ on the same dense open set. This proves all of (C.15).
\end{proof}


\begin{statement}{C4b}{Proved}{An unavoidable failure on a rational homology six-sphere}
Let $f:Y\to\PP^1$ be a connected-fibre holomorphic map from a compact complex threefold with $b_1=b_2=b_3=0$. For every $L\in\Pic(Y)$ there is a scheme-theoretic fibre $F$ with
\[
 H^2(F,(T_Y\otimes L)|_F)\ne0,\qquad
 H^0(F,\Omega_Y^1|_F\otimes(K_Y\otimes L^{-1})|_F)\ne0. \tag{C.16}
\]
For torus general fibre the reduced fibre Euler numbers sum to two. Within the fibre class of C4a, exactly one fibre has rational normalization; it is an irreducible reduced normal-crossing fibre of Euler number two. For that normalization,
\[
H^2(\widetilde S,K_{\widetilde S}\otimes
             \eta^*(K_Y^{-1}\otimes L))=\C
\quad\hbox{for every }L\in\Pic(Y).                    \tag{C.17}
\]
\end{statement}
\begin{proof}
The divisor of any fibre is nonzero and effective. The index argument C.1 therefore gives $H^2(Y,T_Y\otimes M)\ne0$ for every line bundle $M$. Suppose $R^2f_*(T_Y\otimes L)=0$. Twisting by a sufficiently large positive pullback from the base kills $H^1(R^1f_*)$, so Leray would give a line bundle $M=L\otimes f^*\OO(n)$ with $H^2(Y,T_Y\otimes M)=0$, a contradiction. Thus $R^2f_*(T_Y\otimes L)$ is nonzero. Its stalk has nonzero residue at some point, and top-degree base change gives the first nonvanishing in (C.16). Cartier-fibre adjunction and Serre duality give the second, with the exact line bundle displayed there.

For torus general fibre C2a gives $\sum e(F_{\rm red})=e(Y)=2$. In the class of C4a, A.4.6 makes every singular reduced contribution either zero or two, with two exactly for the rational hexagonal normalization; smooth finite torus quotients have zero Euler number. There is consequently exactly one positive fibre. Its smooth rational normalization has $\Pic^0=0$ and torsion-free integral second cohomology. Every ambient line bundle has torsion first Chern class, whose pullback is therefore zero. In particular both $\eta^*K_Y$ and $\eta^*L$ are trivial. Serre duality on the connected rational normalization gives (C.17).
\end{proof}

\begin{statement}{C4c}{Proved}{The finite singular-fibre test and a fixed unavoidable fibre}
Let $f:Y\to\PP^1$ be a proper connected-fibre map from a compact complex threefold with $b_1(Y)=b_2(Y)=0$ and smooth general fibre a complex two-torus. Let $\Sigma$ be its finite set of critical values, and fix a connected component $P$ of $\Pic(Y)$. For $c\in\Sigma$ retain the entire Cartier fibre $F_c$, including its multiplicities, and put
\[
 Z_c=\{L\in P:H^0(F_c,\Omega_Y^1|_{F_c}\otimes
                  (K_Y\otimes L^{-1})|_{F_c})\ne0\}.
                                                               \tag{C.18}
\]
Then the following are equivalent: every $Z_c$ is a proper subset of $P$; and $H^2(Y,T_Y\otimes L)=0$ on a dense open subset of $P$. In that case the generic twist in $P$ satisfies
\[
 H^q(Y,T_Y\otimes L)=0\quad(q\ne1),\qquad
 \chi(Y,T_Y\otimes L)\le0,\qquad c_3(Y)\le0.          \tag{C.19}
\]
For a reduced fibre, membership in the complement of $Z_c$ is exactly the two explicitly typed kernel-and-connecting-map conditions of A.1.1, with kernel $N^*\otimes\mathcal J^{-1}$ computed in A.2. For $F_c=mS$, the equivalent complete test is the finite lifting filtration of A.3.1. These tests use only sheaves and maps on the finitely many singular full fibres.

When $b_3(Y)=0$, at least one fixed critical fibre $F_c$ satisfies $Z_c=P$. Its support is irreducible, by A.3.5. Under the integral assumption $H^2(Y,\Z)=0$, the Picard group is connected, so this is a single fixed fibre for which the ambient group in (C.18) is nonzero for every line bundle on $Y$.
\end{statement}
\begin{proof}
The universal Picard family constructed in C.1 restricts to a holomorphic family on the proper space $F_c\times P$. Its underlying sheaf is flat over $P$, because it is a line-bundle twist of the fixed locally free restriction $\Omega_Y^1|_{F_c}$. Proper semicontinuity therefore makes $Z_c$ a closed complex analytic subset of $P$.

Assume that each $Z_c$ is proper. The complement of their finite union is open dense and has full Lebesgue measure in $P\simeq\C^q$. Statement D.1.11 supplies a dense set of full measure of line bundles non-torsion on every full fibre. Choose $L$ in the intersection. On every smooth torus fibre, $K_Y|_F$ and the conormal are trivial. The intrinsic cotangent bundle is trivial as well, so the conormal exact sequence expresses the ambient cotangent restriction as an extension of trivial bundles. The line $L|_F$ has zero integral Chern class, since $c_1(L)$ is torsion and $H^2(F,\Z)$ is torsion-free; it is nontrivial by the chosen non-torsion condition. Its cohomology in every degree vanishes, as follows by the unitary Fourier resolution of a topologically trivial line on a torus. The ambient section group is consequently zero at every smooth fibre. At each critical fibre it is zero by the definition of the chosen complement.

Serre duality on each Gorenstein Cartier fibre, with $\omega_{F_c}=K_Y|_{F_c}$, and top-degree base change give $R^2f_*(T_Y\otimes L)=0$. Twisting by $f^*\OO(n)$ preserves every fibre restriction. For sufficiently large $n$, it kills $H^1(\PP^1,R^1f_*(T_Y\otimes L)\otimes\OO(n))$. We choose $n$ divisible by the finite order of $c_1(f^*\OO(1))$, so this twist stays in $P$. The two-column Leray sequence gives an actual line bundle in $P$ with zero $H^2$. Proper semicontinuity makes this zero locus dense open. Intersecting it with the dense open endpoint-vanishing loci of C.1 proves (C.19), using $\chi=c_3/2$.

Conversely, suppose $H^2$ vanishes on a dense open subset of $P$, and choose $L$ in that subset and in the full-measure set of D.1.11. The preceding smooth-fibre argument makes $R^2f_*(T_Y\otimes L)$ a torsion sheaf. Leray gives a surjection
\[
 H^2(Y,T_Y\otimes L)\longrightarrow
 H^0(\PP^1,R^2f_*(T_Y\otimes L)).
\]
Thus that torsion sheaf is zero. Top-degree base change and fibre duality imply $L\notin Z_c$ for every $c$. In particular none of the $Z_c$ equals $P$, proving the equivalence.

When $b_3=0$, C.1 gives $H^2(Y,T_Y\otimes L)\ne0$ for every $L$. The equivalence therefore forces some $Z_c$ to be all of $P$. There are only finitely many $c$, and $P$ is an irreducible vector-space torsor; a finite union of proper closed analytic subsets cannot exhaust it. The integral hypothesis makes $P=\Pic(Y)$ by the exponential sequence. This proves the assertion about a fixed fibre. Finally the identifications in A.1--A.3 prove the claimed explicit forms of the fibre tests, without replacing a thickened fibre by its support.
\end{proof}

C4a is an explicit geometric instance of C4c. The two vanishing groups on the support in B1--B2 do not give its fibre test: their ambient space is one-dimensional. The compact quotient A.3.2 proves that the full multiple fibre cannot be replaced by one unshifted reduced-support test. C4c supplies the fixed full-fibre obstruction in arbitrary torus fibrations. C4e proves that its normalization target is nonzero for every twist, with arbitrary singularities and multiplicities.

\begin{statement}{C4d}{Proved}{A finite transgression test for every CDP general-fibre type}
Let $f:Y\to\PP^1$ satisfy the full hypotheses of CDP20 Theorem~2.2: $Y$ is a compact connected complex threefold, $H^1(Y,\Z)=H^2(Y,\Z)=0$, and $f$ is its holomorphic algebraic reduction with connected fibres. Let $\beta\in H^1(Y,\OO_Y)$ lift the integral transgression generator of D.1.9, and put $M_z=\exp(2\pi i z\beta)$, with the same exponential convention as in that statement. For each critical value $c$, define
\[
 B_c=\{z\in\C:
 H^0(F_c,\Omega_Y^1|_{F_c}\otimes M_z|_{F_c})\ne0\},
 \qquad F_c=f^{-1}(c).
 \tag{C.20}
\]
When each $B_c$ is a proper subset of $\C$, a dense open subset of $\Pic(Y)$ satisfies the vanishing and the two inequalities in (C.19). The premise is a finite collection of tests on the actual full fibres, given by A.1--A.3; it is satisfied in particular when every reduced support is smooth.

When $c_3(Y)>0$, there is a fixed critical value $c$ with $B_c=\C$. Its reduced support is nonnormal. Thus a nonnormal reduced fibre is necessary for positive Euler characteristic under the full CDP hypotheses, with arbitrary multiplicities and arbitrary singularities allowed at the other fibres.
\end{statement}
\begin{proof}
The transgression lift exists because the coherent Leray sequence over $\PP^1$ has no degree-one or degree-two base term for $\OO_{\PP^1}$. The integral degree-one Leray sequence identifies the integral invariant group with $\Z$ under the two stipulated integral vanishings. D.1.9 proves that its restriction to every full fibre, and to every nearby smooth fibre, is nonzero. On a neighbourhood $U$ of any fibre, $M_r$ for real $r$ has the unitary local system
\[
 a\longmapsto\exp(2\pi i r\langle\gamma_U,a\rangle)
\]
as its flat realization. This equality follows from the common integral lift $\gamma_U$ in the exponential sequence; no global flat realization on $Y$ is used.

The holomorphic family of exponentiated Cech cocycles representing $z\beta$ restricts to a flat family of coherent bundles on $F_c\times\C$. Proper semicontinuity makes $B_c$ a closed analytic subset of the complex line. A proper such subset is discrete and countable. There are finitely many critical values. We can therefore choose a real irrational $r$ outside the union of all $B_c$ when all are proper.

On a smooth fibre $F$, the character just displayed is the irrational exponential of a nonzero integral class. The explicit calculations for tori, bielliptic surfaces, Kodaira surfaces, Hopf surfaces and Inoue surfaces in A.4.9 give
\[
 H^0(F,\mathbb A_r)=H^1(F,\mathbb A_r)=0.
\]
The unitary maximum principle gives $H^0(F,M_r)=0$, and the injection A.4.7 gives $H^0(F,\Omega_F^1\otimes M_r)=0$. The conormal sequence of a smooth full fibre, whose conormal line is trivial, then gives
$H^0(F,\Omega_Y^1|_F\otimes M_r)=0$.
At every critical full fibre this group vanishes by our choice of $r$. Set $L=K_Y\otimes M_r^{-1}$. Cartier-fibre duality and top-degree base change give $R^2f_*(T_Y\otimes L)=0$. A sufficiently positive base twist kills $H^1(R^1f_*(T_Y\otimes L))$, so
\[
 H^2(Y,T_Y\otimes L\otimes f^*\OO(n))=0
\]
for an actual integer $n$. The exponential sequence identifies $\Pic(Y)$ with a connected complex vector space. Semicontinuity therefore makes degree-two vanishing dense open in this whole group. Intersecting with the endpoint-vanishing loci of C.1 gives (C.19).

For a full fibre with smooth reduction, A.3.4 proves the ambient vanishing for every real irrational $r$, including all the nilpotent layers. Its set $B_c$ is consequently proper. This proves the asserted geometric sufficient case.

Finally $c_3(Y)=2-b_3(Y)$, and the index formula is integral. Thus $c_3(Y)>0$ implies $b_3(Y)=0$. C.1 then gives $H^2(Y,T_Y\otimes L)\ne0$ for every line bundle $L$. The preceding proof shows that the finitely many $B_c$ cannot all be proper; hence one fixed $B_c$ is the entire complex line. Its reduced support cannot be smooth by A.3.4. A connected normal reduced support would be irreducible. Since it is a Cartier divisor in $Y$, its dualizing line has rational first Chern class zero and its structure-sheaf Euler characteristic is zero by divisor Riemann--Roch, exactly as calculated in A.5.2. That statement proves it smooth. We have therefore obtained a nonnormal support, as claimed.
\end{proof}

The fixed fibre in C4d carries an ambient section along the entire complex transgression line. For torus general fibre, C4c strengthens this to all twists in the Picard group. For a general fibre with $h^1(\OO)=1$ and $h^2(\OO)=0$, A.5.13 identifies that line with the entire Picard group. The following proof treats both cases through their actual conductor schemes.

\begin{statement}{C4e}{Proved}{The unrestricted normalization target under the integral CDP hypotheses}
Let $f:Y\to\PP^1$ satisfy the integral CDP hypotheses and $c_3(Y)>0$. Then a fixed irreducible nonnormal reduced support $S$ has normalization $\eta:\widetilde S\to S$ satisfying
\[
 H^2\bigl(\widetilde S,\omega_{\widetilde S}\otimes
           \eta^*(K_Y^{-1}\otimes L)\bigr)\ne0
 \qquad\hbox{for every }L\in\Pic(Y).                  \tag{C.21}
\]
The assertion concerns the reduced support; no multiplicity-one conclusion is included.
\end{statement}
\begin{proof}
C.3.1 and the exclusions C.3.7, C.3.9 and C.3.10 show
that the general fibre is a two-torus, a Hopf surface, or
an Inoue surface. In the last two cases it has
$h^1(\OO)=1$ and $h^2(\OO)=0$.
For torus general fibre, C4c gives a fixed full fibre with a nonzero ambient cotangent section for every ambient twist. In the second case, C4d gives persistence along $\C\beta$, and A.5.13 identifies this with $\Pic(Y)$. From this point the proof uses precisely that persistence property and does not use $c_3(Y)>0$ again. A.3.5 makes its support irreducible, so it is $mS$ for a single positive integer $m$. A.3.4 excludes a smooth support, and A.5.2 then excludes a normal support. Its filtration in A.3.1 has finitely many graded factors
\[
 \Omega_Y^1|_S\otimes M|_S\otimes\OO_Y(-jS)|_S,
 \qquad 0\le j<m.
\]
Nonvanishing on the filtered sheaf requires nonvanishing on at least one factor. Each factor's nonvanishing locus is closed analytic in the connected vector group $\Pic(Y)$. A finite union of proper such loci cannot cover that group. Hence one factor has nonzero sections for every $M$. Translating by $\OO_Y(-jS)$ shows
$H^0(S,\Omega_Y^1|_S\otimes M|_S)\ne0$ for every $M$.
A.4.9 kills the torsion-free quotient on a dense open set of twists. The exact sequence A.1.1 and proper semicontinuity then show that the inverse-Jacobian kernel itself has sections for every twist.

We apply the precise persistent-kernel reduction A.5.4. It gives a fixed component of the normalization on whose resolution every twist of one fixed rank-one sheaf has a section. The Gauduchon degree argument there proves that the pullback of every ambient line has degree zero for every Gauduchon metric. Its conductor comparison excludes nonnegative Kodaira dimension, and its class VII argument makes every ambient line trivial on that normalization component. Rational components have the same triviality because their Picard group has no connected part. Either case gives (C.21) immediately by Serre duality and descent from resolution. It remains to treat an irrational ruled resolution $p:Z\to B$.

Let $h:Z\to S\subset Y$. A.5.3 supplies an effective conductor divisor $D_Z$ with
\[
 K_Z+D_Z=h^*(K_Y+S),\qquad D_Z\cdot F=2              \tag{C.22}
\]
for a general ruling fibre $F$. Its horizontal part is therefore either reduced of total ruling degree two, or $2C$ for a section $C$. We record the comparison through each ruling blowdown $\sigma:Z\to Z'$ used below. Put $D'=\sigma_*D_Z$, let $a$ be the coefficient of its exceptional curve $E$ in $D_Z$, and put $d=\operatorname{mult}_{\sigma(E)}D'$. Then
\[
 K_Z+D_Z=\sigma^*(K_{Z'}+D')+(1+a-d)E.
 \tag{C.22b}
\]
Intersection with $E$, using $c_1(K_Z+D_Z)=0$ and $E^2=-1$, gives $a=d-1$. Thus the log-canonical divisor is the exact pullback of $K_{Z'}+D'$, and the latter again has zero rational Chern class. No factorization of $h$ through the blowdown is asserted or needed.

We first treat the reduced horizontal part $H$. If a component is exceptional for the normalization resolution, every ambient line restricts trivially to it. Its normalization maps finitely onto $B$, and the norm composed with pullback on $\Pic^0(B)$ is multiplication by the covering degree. Hence this restriction has finite kernel; the connected ambient Picard image is trivial, giving (C.21). We may therefore assume that the horizontal components map to conductor curves. Adjunction on each of them gives
\[
 2p_a(H_i)-2+(D_Z-H_i)\cdot H_i=0.
\]
The second term is nonnegative. Since $H_i$ dominates a positive-genus base, its arithmetic genus is at least one. Thus the base has genus one, each $H_i$ is smooth elliptic, and the horizontal curves are disjoint from each other and from the vertical part. Contract ruling $(-1)$-curves to a minimal bundle. A contracted curve outside $D_Z$ has $D_Z\cdot E=1$, so cannot identify two horizontal points or create a horizontal node. A contracted curve in the vertical part is disjoint from $H$. Downstairs the horizontal curves remain smooth and disjoint, and any nonzero vertical divisor on a minimal ruled bundle meets them. Consequently the downstairs vertical part is zero. Reversing the contractions, a blowup away from the horizontal conductor would give coefficient $-1$ on its exceptional curve in (C.22), which persists under further strict transform and contradicts effectiveness. A blowup at a smooth point of a coefficient-one horizontal curve gives exceptional coefficient zero. Hence $D_Z=H$ has no vertical component. Every codimension-one singularity of $S$ is then an ordinary node, and A.2.3 identifies the inverse-Jacobian kernel with the normalization kernel globally, including its codimension-two points. A section of a topologically trivial line on the Kahler resolution has an effective zero divisor of zero Kahler volume, so is nowhere zero. Persistence for all twists therefore makes all ambient pullbacks trivial and proves (C.21).

We now have $D_Z=2C+V$, where $C$ is a section and $V$ is vertical. If $C$ is exceptional for $Z\to\widetilde S$, every ambient line is trivial on $C$. Restriction $\Pic^0(Z)\to\Pic^0(C)$ is an isomorphism, because $C$ is a ruling section. This again makes all ambient pullbacks trivial. Otherwise let $D\subset S$ be its image and let $D^\nu$ normalize that curve. When $D^\nu$ is rational, every ambient line pulls back trivially to $D^\nu$, hence to $C$, giving the same conclusion.

Suppose instead that $g(D^\nu)\ge1$. A vertical rational curve cannot map finitely onto $D^\nu$, so all the normalization branches over the general point of $D$ belong to $C$. Their conductor coefficient is two. We derive the possible transverse reduced plane curves rather than assuming a list of singularities. Write its equation as $f=\prod_i f_i$ in the smooth-plane local ring and let $R_i$ and $B_i$ be a branch ring and its normalization. The elements of $R=\OO/(f)$ supported only on branch $i$ form
\[
 \left(\prod_{j\ne i}f_j\right)\bigg|_{f_i}\,R_i.
\]
This equality follows from
$\bigcap_{j\ne i}(f_j)=(\prod_{j\ne i}f_j)$ in that unique
factorization domain. Multiplying by arbitrary elements of
$B_i$ therefore shows that the conductor exponent on branch $i$
is its intrinsic exponent plus $\sum_{j\ne i}I(f_i,f_j)$.
A singular branch has intrinsic exponent at least two:
exponent one would imply $\C+t\C\{t\}=\C\{t\}\subset R_i$.
Exponent two forces $R_i=\C+t^2\C\{t\}=\C\{t^2,t^3\}$,
since an element of order one would again make the ring smooth.
Any additional branch has intersection multiplicity at least
two with a singular branch and is therefore excluded. Thus a
singular branch gives exactly the cusp. When all branches are
smooth, their pairwise positive intersection multiplicities sum
to two on each branch. This gives two branches of contact two
(the tacnode), or three pairwise transverse branches (the
ordinary triple point), and exhausts the possibilities.
The only remaining alternative is a nonconstant pullback
$\Pic(Y)\to\Pic(Z)$; a constant pullback is the zero
homomorphism and already gives (C.21). In that alternative,
the coherent conductor argument A.5.12 gives $p_a(D)\le1$
for torus general fibre, and A.5.13 gives the same bound
in the second case.
Since $g(D^\nu)\ge1$, it follows that $D$ itself is smooth
elliptic. This conclusion permits nonreduced conductor image
schemes and imposes no assumption on the original singular
curve. We can therefore apply C2d directly. It makes the
transverse type constant, gives the etale conductor cover
$C\to D$, and proves $C^2=0$. In particular $C$ and the
ruling base are elliptic.

Adjunction in (C.22) gives $V\cdot C=0$. No ruling $(-1)$-curve meets $C$: for one outside $V$, its intersection with $K_Z+2C+V$ would be at least $-1+2=1$; for one in $V$, the equality $V\cdot C=0$ forbids such an intersection. Thus every contraction to the minimal ruled bundle avoids $C$ and preserves its self-intersection zero. On that bundle $K+2C$ is numerically trivial, so the effective vertical part is numerically trivial and hence zero. A first reverse blowup, necessarily away from $C$, would again produce negative conductor coefficient. There are therefore no reverse blowups: $Z$ is minimal ruled and $D_Z=2C$. There is also no exceptional curve for the normalization resolution: a negative curve $E\ne C$ would satisfy $2p_a(E)-2=E^2-2C\cdot E<0$, hence would be rational and contained in a ruling fibre; a minimal ruled bundle has only self-intersection-zero fibres. Thus the normalization itself is smooth here.

By the explicit inverse-Jacobian computations in C2d, its pullback allows at most a simple pole along $C$. There are no other divisorial poles. Hartogs extension for the line bundle on the smooth surface gives an injection of the kernel into $h^*\OO_Y(-S)\otimes\OO_Z(C)$. The exact sequence of the section $C$ and $R^1p_*\OO_Z=0$ give
\[
 0\longrightarrow\OO_B\longrightarrow p_*\OO_Z(C)
   \longrightarrow N_{C/Z}\longrightarrow0,
 \qquad \deg N_{C/Z}=0.                              \tag{C.22a}
\]
For a degree-zero line $P$ on the elliptic curve, both end terms after tensoring by $P$ have zero sections unless $P=\OO_B$ or $P=N_{C/Z}^{-1}$. Thus $H^0(Z,\OO_Z(C)\otimes p^*P)$ is supported at at most two points of $\Pic^0(B)$. The ambient pullback image is a connected complex subgroup of this Picard group. Persistence for every ambient twist forces that subgroup to be a point, and it contains the identity. Therefore every ambient line pulls back trivially to $Z$ in this last case as well.

The equality $\pi_*\OO_Z=\OO_{\widetilde S}$ for the normal resolution descends triviality of a pulled-back line: projection formula identifies its space of sections with that of its trivial pullback, and the same statement for its inverse gives mutually inverse trivializations. Consequently one normalization component has trivial $\eta^*(K_Y^{-1}\otimes L)$ for every $L$. Serre duality on that compact normal Cohen--Macaulay surface identifies its top dualizing cohomology with $H^0(\OO)^\vee=\C$. It is a direct summand of the target in (C.21), proving the claim.
\end{proof}

The local example A.2.5 has actual vertical inverse-Jacobian
poles, and A.2.6 has a critical curve that is not a local
complete intersection. Their local rings remain part of the
audit. The global conductor-duality arguments A.5.12--A.5.13 prove
that a positive-genus conductor image in the persistent ruled
case is smooth, and the exact contractions above then remove
the vertical conductor. Thus neither local phenomenon leaves
an unproved step in the stated theorem. The assertion
is about the normalization of a reduced support and does not
claim that its full fibre has multiplicity one.

\begin{statement}{C4f}{Proved}{An unrestricted geometric replacement for CDP20 Theorem 2.2}
Let $f:Y\to\PP^1$ satisfy the integral CDP hypotheses.
For each critical full fibre with
irreducible nonnormal reduced support $S$, let
$\eta_S:T_S\to S$ be its normalization. Then a general
line bundle $L\in\Pic(Y)$ satisfies
\[
 H^q(Y,T_Y\otimes L)=0\quad(q\ne1)
 \tag{C.G.1}
\]
exactly when, for every such $S$, the holomorphic pullback
homomorphism
\[
 \Pic(Y)\longrightarrow\Pic(T_S)
 \tag{C.G.2}
\]
is nonzero. Equivalently, its differential
$H^1(Y,\OO_Y)\to H^1(T_S,\OO_{T_S})$ is nonzero.
In this case $\chi(Y,T_Y\otimes L)\le0$ and $c_3(Y)\le0$.
If one map (C.G.2) is zero, then
$H^2(Y,T_Y\otimes L)\ne0$ for every $L\in\Pic(Y)$.
All fibre multiplicities and all reduced-support
singularities are permitted in this criterion.
\end{statement}
\begin{proof}
The classification and exclusions cited in C4e leave
exactly the two cohomological cases used in that proof:
torus general fibre, or $h^1(\OO)=1$, $h^2(\OO)=0$.
The Picard group of $Y$ is the connected exponential image
of $H^1(Y,\OO_Y)$. A holomorphic group homomorphism out of
this group is zero precisely when its differential is zero.
For each fixed critical full fibre, the ambient-section
nonvanishing locus is closed analytic in this group, as
in C4c. A.3.4 makes it proper when the support is smooth,
and A.3.5 makes it proper when the support is reducible.
A connected normal support is irreducible and is smooth
by A.5.2, so these cases leave just the supports in the
statement.

Suppose the nonvanishing locus for $F=mS$ is the whole
Picard group. The proof of C4e, beginning with its
fixed-fibre persistence assertion, proves that (C.G.2)
is zero. Its steps use the actual component filtration,
the inverse-Jacobian kernel, the Gauduchon bound, the
effective conductor comparison, the coherent conductor
genus bound, and the explicit constant-transverse-type
calculations. Each is independent of the Euler number
once persistence has been given.

Conversely, suppose (C.G.2) is zero. The last nonzero
ideal layer of the full fibre is an actual subsheaf
\[
 (\Omega_Y^1\otimes M\otimes\OO_Y(-(m-1)S))|_S
 \hookrightarrow
          (\Omega_Y^1\otimes M)|_{mS}.
 \tag{C.G.3}
\]
The canonical normalization submodule in A.2.2 gives
an injection into the left side from
\[
 \eta_{S*}\bigl(\eta_S^*((M\otimes\OO_Y(-mS))|_S)\bigr).
 \tag{C.G.4}
\]
Here the first conormal factor $\OO_Y(-S)|_S$ combines
with the displayed $(m-1)$st ideal layer; this comparison
retains every nilpotent layer needed for the injection.
Since $mS$ is a full fibre, $\OO_Y(-mS)|_S$ is a trivial
line, with its one-dimensional scalar factor determined
by a local base parameter. The pullback of $M$ to $T_S$
is trivial by (C.G.2). Hence (C.G.4) has a nonzero global
section for every $M$, and its injections prove that the
full-fibre nonvanishing locus is the entire Picard group.
This also proves the converse directly for arbitrary $m$.

We have proved an exact geometric characterization of
every improper critical-fibre locus. For torus general
fibre, applying C4c gives the equivalence with generic
degree-two vanishing. In the second case A.5.13 identifies
$\Pic(Y)$ with the transgression line, so properness of
all these loci is exactly the premise of C4d and gives
generic degree-two vanishing there as well.
The endpoint vanishings of C.1 then give (C.G.1) and the
index inequalities.
For the last assertion in the torus case, generic non-torsion on every
smooth torus fibre makes $R^2f_*(T_Y\otimes L)$ a
torsion sheaf for a dense set of twists, by D.1.11.
The fixed full-fibre obstruction makes it nonzero.
Its nonzero global sections are a quotient of
$H^2(Y,T_Y\otimes L)$ by the Leray sequence.
Thus degree-two cohomology is nonzero on a dense set.
Its proper semicontinuity locus is closed analytic in
the connected Picard family, so it is nonzero for every
twist. In the second case use instead the actual twists
$L_r=K_Y\otimes M_r^{-1}$ for real irrational $r$.
A.4.9 and the smooth-fibre conormal sequence make their
$R^2f_*(T_Y\otimes L_r)$ torsion. The fixed full-fibre
obstruction again makes that torsion sheaf nonzero,
so Leray gives $H^2(Y,T_Y\otimes L_r)\ne0$.
These twists have accumulation points in the complex
line $\Pic(Y)=\C\beta$. A closed complex analytic subset
of that line containing them is the whole line.
This proves all-twist nonvanishing in the second case
and completes both implications of the criterion.
\end{proof}


\begin{statement}{C4g}{Proved}{The necessary normalization defect in the three named test families}
For every admissible member of the original $(3,4)$ family,
every member of the two-parameter family in E13, and every
member of the Liu--Rollenske family in E14, there is a reduced
full hexagonal fibre $W$ with normalization
$\eta:\mathrm dP_6\to W$ such that
\[
 \eta^*L=\OO_{\mathrm dP_6}\quad(L\in\Pic(Y)),\qquad
 H^2\bigl(\mathrm dP_6,K_{\mathrm dP_6}\otimes
                    \eta^*(K_Y^{-1}\otimes L)\bigr)=\C.
 \tag{C.TEST.1}
\]
In these three families the singular-fibre Euler contributions are
\[
\begin{array}{c|c|c}
\text{family}&\text{reduced singular supports}&
                 \text{their Euler numbers}\\ \hline
(3,4)&W&2\\
\text{two-parameter E13}&W,\ W_1,\ W_2&2,0,0\\
\text{Liu--Rollenske E14}&W,\ C_{\rm cusp}&2,0 .
\end{array}
\tag{C.TEST.2}
\]
The other critical full fibres are smooth multiple fibres,
with multiplicities $(3,4)$, $4$, and $4$, respectively.
Every one of those full multiple fibres passes the generic
ambient vanishing test A.3.4, including all its nilpotent
layers. Consequently positive $c_3$ in these examples does
not require a failure of that test on a nonreduced fibre.
The reduced full fibre $W$ already supplies the fixed
normalization obstruction, and $H^2(Y,T_Y\otimes L)\ne0$
for every ambient line bundle in each family.
\end{statement}
\begin{proof}
The original local filling is the hexagonal toric quotient
in A.4.1. The two-parameter construction in E13 places this
same filling at the simple nodal value where its specified
linearization vanishes; its other two simple nodal values
have rank-one fillings. The coordinate and line-bundle
comparisons of E14 place the hexagonal filling at its
separate simple node and the cuspidal quotient at the
remaining critical value. Those proofs construct the actual
local spaces and their identifications over the punctured
base discs. Thus the normalization of $W$ in each case is
the stated smooth rational surface.

These threefolds have $H^2(Y,\Z)=0$. Pulling back any
ambient line to $\mathrm dP_6$ gives zero integral first
Chern class. The exponential sequence there, with
$H^1(\mathrm dP_6,\OO)=0$, makes that line trivial.
This applies to $K_Y^{-1}\otimes L$ as well as to $L$.
Serre duality on $\mathrm dP_6$ proves (C.TEST.1).
The last assertion about ambient $H^2$ follows from C4f
through the zero Picard pullback at this fixed fibre.
It does not require surjectivity of the global
normalization map for every twist; C5b computes the
actual exceptions to that surjectivity in the original
family.

The toric strata of $W$ consist of one $(\C^*)^2$ stratum,
three $\C^*$ strata and two triple points. Their compactly
supported Euler numbers give $e(W)=2$.
Each rank-one normalization is a $\PP^1$-bundle over an
elliptic curve, with its zero and infinity sections glued.
Both the normalization and these conductor curves have
Euler number zero, so normalization additivity gives
$e(W_i)=0$. At the cusp in E14, normalization of the
cuspidal rational curve makes its degree-zero line bundle
trivial. The lifted contracting action on this normalized
line complement is a product of a fixed automorphism of
$\PP^1$ and multiplication by a constant of modulus
strictly less than one on $\C^*$. Its quotient is a
$\PP^1$-bundle over the elliptic curve given by that
multiplicative period. The cusp normalization is
bijective; a proper continuous bijection between these
compact Hausdorff spaces is a homeomorphism. Hence
$e(C_{\rm cusp})=0$. This proves (C.TEST.2).

Finally the explicit good-reduction covers in E13--E14
and the original root charts have the stated orders.
Their reduced central fibres are smooth finite unramified
torus quotients, and their Euler numbers are zero.
A.3.4 applies to their actual full fibres, proving the
asserted generic ambient vanishing. Euler localization
therefore gives $c_3(Y)=2$ in all three rows, while the
nonreduced fibres satisfy the complete multiple-fibre
test. This is the required test of the proposed additional
nonreduced-fibre necessity.
\end{proof}

\paragraph{Literature comparison for the geometric replacement.}
The ambient-kernel failure for the original construction is
located in \cite[Section~10]{AC26}; the general kernel is
calculated in A.1--A.2. Statements C4e--C4f use the actual
conductor-duality map and the independently proved
smooth-fibre exclusions to obtain a geometric criterion
under all the integral hypotheses of \cite[Theorem~2.2]{CDP20}.
The latter paper already states the Kodaira and bielliptic
exclusions in its Corollary~5.3, with its Corollary~5.2 as
an input; C.3.7 and C.3.9--C.3.10 prove those exclusions
without that input. The Inoue exclusion C.3.12 is derived
from the explicit complete families of \cite{Inoue1974}.
Searches on 15 November 2026 for compact threefolds,
vanishing second Betti number, Kodaira or Inoue general
fibres, normalization obstructions, and degenerate-torus
Euler numbers did not supply the all-singularity
criterion C4f in the sources opened. We claim no priority.
